% ===== latex-rendered:document =====
% flm: latex-rendered; style file flmx.5f662687.sty
\documentclass{article}
\usepackage{flmx.5f662687}
\begin{document}
The claim holds in every case.\flmxSpan[style=superscript, role=flm.xnotes.xnote.mark, anchor=fnref--1yuhk]{\flmxLink[kind=xnote, href={\#fn-1}, target=fn-1, doc-role=noteref, label={Footnote 1}]{1}}
The next sentence goes on.

Two paragraphs.\flmxSpan[style=superscript, role=flm.xnotes.xnote.mark, anchor=fnref--2yurg]{\flmxLink[kind=xnote, href={\#fn-2}, target=fn-2, doc-role=noteref, label={Footnote 2}]{2}}
And on.

\begin{flmxBlock}[role=flm.xnotes.footnote.list, kinds={xnotes footnote}, doc-role=endnotes]
\flmxHeading[kind=xnotes, level=2, toc=false, plain={Footnotes}]{Footnotes}
\begin{flmxList}[kinds={numbered xnotes footnote labeled}, depth=1]
\flmxItem[label={\flmxSpan[style=superscript, role=flm.xnotes.entry.mark]{1}}, anchor=fn-1]
Except the empty one, which we exclude by convention.

\flmxItem[label={\flmxSpan[style=superscript, role=flm.xnotes.entry.mark]{2}}, anchor=fn-2]
First paragraph.

Second paragraph.
\end{flmxList}
\end{flmxBlock}

\end{document}
% ===== latex-rendered:style flmx.5f662687.sty =====
\NeedsTeXFormat{LaTeX2e}[2022/06/01]
\ProvidesPackage{flmx.5f662687}
% the definitions of a document written by FLM's latex-rendered backend
% the engine: LuaLaTeX
\RequirePackage{iftex}
\RequireLuaTeX
% the packages
% flm.latex.pkg.fontspec
\RequirePackage{fontspec}
% flm.latex.pkg.amsmath
\RequirePackage{amsmath}
% flm.latex.pkg.xcolor
\RequirePackage{xcolor}
% hyperref, loaded at the end of the preamble
\AddToHook{begindocument/before}{\RequirePackage[colorlinks=true, linkcolor=flmxLinkColor, urlcolor=flmxUrlColor, citecolor=flmxCiteColor]{hyperref}}
% flm.latex.def.flmx.core
% flm.latex.def.flmx.core -- the core of the latex-rendered contract.
%
% The rendered LaTeX writes one command or environment per presentation
% element, all named \flmx...; this chunk defines the machinery every one of
% them uses, the paragraph-level elements (flmxBlock, \flmxSpan), the three
% decorations, the target of an anchor, two lengths and the colors.
%
% Keys. Every element takes an optional key-value argument first, holding
% the facts FLM knows about it ([role=..., kinds={a b}, anchor=...]). While
% an element is written, its keys are in scope:
%   \flmxKey{<key>}               expands to the key's value, empty where the
%                                 element does not carry the key;
%   \flmxIfKey{<key>}{<yes>}{<no>} tests whether the element carries the key.
% Each element sees its own keys only: a block inside a block sees its own,
% and the outer block's are back after it.
%
% Registering a look. A template changes how an element looks by registering
% code under a role, a kind token, a style name, a heading kind or level, or a
% link kind; a later registration wins:
%   \flmxDefineBlock{<role or kind token>}{<begin code>}{<end code>}
%   \flmxDefineSpan{<role, kind token or style name>}{<code over #1>}
%   \flmxDefineHeading{<kind name or level>}{<code over #1>}
%   \flmxDefineRuninHeading{<kind name or level>}{<code over #1>}
%   \flmxDefineLink{<kind>}{<code over #1>}
%   \flmxDefineList{<kind token>}{<begin code>}{<end code>}
%   \flmxDefineItem{<kind token>}{<code run at the item>}
%   \flmxDefineVerbatim{<kind>}{<code over #1>}
%   \flmxDefineVerbatimBlock{<kind>}{<begin code>}{<end code>}
%   \flmxDefineXnote{<kind>}{<code over #1 (the mark) and #2 (the content)>}
%   \flmxDefineCellStyle{<style name>}{<tblr cell keys>}
% An element is written with the first registration found, in this order:
%   flmxBlock: its role, then each kind token in order;
%   \flmxSpan: its role, then each kind token, then its style;
%   \flmxHeading, \flmxRuninHeading: its kind, then its level;
%   \flmxLink, \flmxVerbatim, flmxVerbatimBlock, \flmxXnote: its kind;
%   flmxList, \flmxItem: each kind token of the list in order;
%   a cell of flmxCells: the keys registered for each of its styles, in
%   order, added to the cell's own;
% and with the element's default where none is registered. Each element's
% default is replaced with
%   \flmxDefine<Element>Default{<code>}
% (\flmxDefineBlockDefault{<begin>}{<end>}, \flmxDefineSpanDefault{<code
% over #1>}, and so on for Heading, RuninHeading, Link, List, Item, Cells,
% Lines, Line, Verbatim, VerbatimBlock, Image, Xnote, the three decorations
% DecoRunin, DecoFinal, DecoMargin, and the three marks Dropped,
% Unsupported and Unknown, whose code is over #1 and #2).
%
% The elements this chunk defines:
%   \begin{flmxBlock}[role, kinds, anchor, doc-role] ... \end{flmxBlock}
%       a block of paragraphs; default: a paragraph break before and after,
%       the target of its anchor first.
%   \flmxSpan[style, role, kinds, anchor, doc-role]{<content>}
%       inline content; defaults by style: emph \emph, bold \textbf, italic
%       \textit, monospace \texttt, underline \underline, small-caps \textsc,
%       superscript \textsuperscript, subscript \textsubscript,
%       unresolved-reference in flmxMarkColor; no style: the content.
%   \flmxDecoRunin[stop]{<content>}  the title set at the start of a
%       paragraph: the content and the stop in bold, then \flmxRuninSep.
%   \flmxDecoFinal{<content>}  set flush right at the end of the last line;
%       where it stands alone (after a list, for instance), it starts a
%       paragraph of its own.
%   \flmxDecoMargin{<content>}  a margin note, in \footnotesize.
%   \flmxTarget{<anchor>}
%       makes the current place a target named <anchor>: in text,
%       \phantomsection and \label{flmx:<anchor>}; in mathematics, a
%       hyperlink target and the label written by hand. Redefine it to change
%       the scheme; the prefix flmx: of the label is fixed, so that
%       \pageref{flmx:<anchor>} and \hyperref[flmx:<anchor>]{...} find it.
%
% Lengths and colors (redefine with \setlength and \colorlet):
%   \flmxTabWidth   one tab of an indentation, 2em;
%   \flmxRuninSep   the space after a run-in title, 0.8em;
%   flmxLinkColor, flmxUrlColor, flmxCiteColor  the link colors hyperref is
%                   given; flmxMarkColor  the color of what FLM could not
%                   write (the marks and an unresolved reference).
%
% Paragraphs are LaTeX paragraphs; this chunk changes no paragraph setting.

\newlength\flmxTabWidth
\setlength\flmxTabWidth{2em}
\newlength\flmxRuninSep
\setlength\flmxRuninSep{0.8em}
\colorlet{flmxLinkColor}{blue!60!black}
\colorlet{flmxUrlColor}{blue!60!black}
\colorlet{flmxCiteColor}{green!40!black}
\colorlet{flmxMarkColor}{red}

\ExplSyntaxOn

% --- the keys of the element being written ---------------------------------

\seq_new:N \l__flmx_keys_seq
\tl_new:N \l__flmx_kinds_tl
\tl_new:N \l__flmx_found_tl
\bool_new:N \l__flmx_found_bool

\cs_new_protected:Npn \__flmx_set_keys:n #1
  {
    \seq_map_inline:Nn \l__flmx_keys_seq
      { \cs_set_eq:cN { l__flmx_key_##1_tl } \tex_undefined:D }
    \seq_clear:N \l__flmx_keys_seq
    \keyval_parse:nnn { \__flmx_key_flag:n } { \__flmx_key_value:nn } {#1}
  }
\cs_new_protected:Npn \__flmx_key_flag:n #1
  { \__flmx_key_value:nn {#1} { true } }
\cs_new_protected:Npn \__flmx_key_value:nn #1#2
  {
    \cs_set_nopar:cpn { l__flmx_key_#1_tl } {#2}
    \seq_put_right:Nn \l__flmx_keys_seq {#1}
  }

\NewExpandableDocumentCommand \flmxKey { m }
  { \cs_if_exist_use:c { l__flmx_key_#1_tl } }
\NewExpandableDocumentCommand \flmxIfKey { m m m }
  { \cs_if_exist:cTF { l__flmx_key_#1_tl } {#2} {#3} }

% --- finding the registered code ------------------------------------------

% #1 the element, #2 the keys to try in order (`kinds' tries each token).
% Leaves the name of the code to run in \l__flmx_found_tl.
\cs_new_protected:Npn \__flmx_find:nn #1#2
  {
    \bool_set_false:N \l__flmx_found_bool
    \tl_set:Nn \l__flmx_found_tl { __flmx_#1_default }
    \clist_map_inline:nn {#2}
      {
        \str_if_eq:nnTF {##1} { kinds }
          {
            \cs_if_exist:cT { l__flmx_key_kinds_tl }
              {
                \tl_set_eq:Nc \l__flmx_kinds_tl { l__flmx_key_kinds_tl }
                \tl_replace_all:Nnn \l__flmx_kinds_tl { ~ } { , }
                \clist_map_inline:Nn \l__flmx_kinds_tl
                  { \__flmx_try:nn {#1} {####1} }
              }
          }
          {
            \cs_if_exist:cT { l__flmx_key_##1_tl }
              { \__flmx_try:nv {#1} { l__flmx_key_##1_tl } }
          }
      }
  }
\cs_new_protected:Npn \__flmx_try:nn #1#2
  {
    \bool_if:NF \l__flmx_found_bool
      {
        \cs_if_exist:cT { __flmx_#1_reg_#2 }
          {
            \tl_set:Nn \l__flmx_found_tl { __flmx_#1_reg_#2 }
            \bool_set_true:N \l__flmx_found_bool
          }
      }
  }
\cs_generate_variant:Nn \__flmx_try:nn { nv }

% --- writing an element ------------------------------------------------------

% A command of one content argument: #1 element, #2 keys to try, #3 the
% key-value list, #4 the content.
\cs_new_protected:Npn \__flmx_command:nnnn #1#2#3#4
  {
    \group_begin:
      \__flmx_set_keys:n {#3}
      \__flmx_find:nn {#1} {#2}
      \use:c { \l__flmx_found_tl } {#4}
    \group_end:
  }
% A command of two arguments.
\cs_new_protected:Npn \__flmx_command:nnnnn #1#2#3#4#5
  {
    \group_begin:
      \__flmx_set_keys:n {#3}
      \__flmx_find:nn {#1} {#2}
      \use:c { \l__flmx_found_tl } {#4} {#5}
    \group_end:
  }
% The begin and end of an environment (LaTeX's environment is the group).
\cs_new_protected:Npn \__flmx_env_begin:nnn #1#2#3
  {
    \__flmx_set_keys:n {#3}
    \__flmx_find:nn {#1} {#2}
    \cs_set_eq:cN { l__flmx_#1_env_tl } \l__flmx_found_tl
    \use:c { \l__flmx_found_tl _begin }
  }
\cs_new_protected:Npn \__flmx_env_end:n #1
  { \use:c { \use:c { l__flmx_#1_env_tl } _end } }

% Registering code: #1 the element, #2 the name, #3 the code.
\cs_new_protected:Npn \__flmx_register:nnn #1#2#3
  { \cs_set_protected:cpn { __flmx_#1_reg_#2 } ##1 {#3} }
\cs_new_protected:Npn \__flmx_register_env:nnnn #1#2#3#4
  {
    \cs_set_protected:cpn { __flmx_#1_reg_#2 } { }
    \cs_set_protected:cpn { __flmx_#1_reg_#2_begin } {#3}
    \cs_set_protected:cpn { __flmx_#1_reg_#2_end } {#4}
  }
\cs_new_protected:Npn \__flmx_default:nn #1#2
  { \cs_set_protected:cpn { __flmx_#1_default } ##1 {#2} }
\cs_new_protected:Npn \__flmx_default_ii:nn #1#2
  { \cs_set_protected:cpn { __flmx_#1_default } ##1##2 {#2} }
\cs_new_protected:Npn \__flmx_default_none:nn #1#2
  { \cs_set_protected:cpn { __flmx_#1_default } {#2} }
\cs_new_protected:Npn \__flmx_default_env:nnn #1#2#3
  {
    \cs_set_protected:cpn { __flmx_#1_default_begin } {#2}
    \cs_set_protected:cpn { __flmx_#1_default_end } {#3}
  }

% The target of the element's anchor, where it carries one.
\cs_new_protected:Npn \__flmx_anchor:
  { \flmxIfKey { anchor } { \flmxTarget { \flmxKey { anchor } } } { } }

% --- the registrations ---------------------------------------------------

\NewDocumentCommand \flmxDefineBlock { m +m +m }
  { \__flmx_register_env:nnnn { block } {#1} {#2} {#3} }
\NewDocumentCommand \flmxDefineSpan { m +m }
  { \__flmx_register:nnn { span } {#1} {#2} }
\NewDocumentCommand \flmxDefineHeading { m +m }
  { \__flmx_register:nnn { heading } {#1} {#2} }
\NewDocumentCommand \flmxDefineRuninHeading { m +m }
  { \__flmx_register:nnn { runinheading } {#1} {#2} }
\NewDocumentCommand \flmxDefineLink { m +m }
  { \__flmx_register:nnn { link } {#1} {#2} }
\NewDocumentCommand \flmxDefineList { m +m +m }
  { \__flmx_register_env:nnnn { list } {#1} {#2} {#3} }
\NewDocumentCommand \flmxDefineItem { m +m }
  { \cs_set_protected:cpn { __flmx_item_reg_#1 } {#2} }
\NewDocumentCommand \flmxDefineVerbatim { m +m }
  { \__flmx_register:nnn { verbatim } {#1} {#2} }
\NewDocumentCommand \flmxDefineVerbatimBlock { m +m +m }
  { \__flmx_register_env:nnnn { verbatimblock } {#1} {#2} {#3} }
\NewDocumentCommand \flmxDefineXnote { m +m }
  { \cs_set_protected:cpn { __flmx_xnote_reg_#1 } ##1##2 {#2} }
\NewDocumentCommand \flmxDefineCellStyle { m +m }
  {
    \tl_clear_new:c { l__flmx_cellstyle_#1_tl }
    \tl_set:cn { l__flmx_cellstyle_#1_tl } {#2}
  }

\NewDocumentCommand \flmxDefineBlockDefault { +m +m }
  { \__flmx_default_env:nnn { block } {#1} {#2} }
\NewDocumentCommand \flmxDefineSpanDefault { +m }
  { \__flmx_default:nn { span } {#1} }
\NewDocumentCommand \flmxDefineHeadingDefault { +m }
  { \__flmx_default:nn { heading } {#1} }
\NewDocumentCommand \flmxDefineRuninHeadingDefault { +m }
  { \__flmx_default:nn { runinheading } {#1} }
\NewDocumentCommand \flmxDefineLinkDefault { +m }
  { \__flmx_default:nn { link } {#1} }
\NewDocumentCommand \flmxDefineListDefault { +m +m }
  { \__flmx_default_env:nnn { list } {#1} {#2} }
\NewDocumentCommand \flmxDefineItemDefault { +m }
  { \__flmx_default_none:nn { item } {#1} }
\NewDocumentCommand \flmxDefineCellsDefault { +m +m }
  { \__flmx_default_env:nnn { cells } {#1} {#2} }
\NewDocumentCommand \flmxDefineLinesDefault { +m +m }
  { \__flmx_default_env:nnn { lines } {#1} {#2} }
\NewDocumentCommand \flmxDefineLineDefault { +m }
  { \__flmx_default:nn { line } {#1} }
\NewDocumentCommand \flmxDefineVerbatimDefault { +m }
  { \__flmx_default:nn { verbatim } {#1} }
\NewDocumentCommand \flmxDefineVerbatimBlockDefault { +m +m }
  { \__flmx_default_env:nnn { verbatimblock } {#1} {#2} }
\NewDocumentCommand \flmxDefineImageDefault { +m }
  { \__flmx_default:nn { image } {#1} }
\NewDocumentCommand \flmxDefineXnoteDefault { +m }
  { \__flmx_default_ii:nn { xnote } {#1} }
\NewDocumentCommand \flmxDefineDecoRuninDefault { +m }
  { \__flmx_default:nn { decorunin } {#1} }
\NewDocumentCommand \flmxDefineDecoFinalDefault { +m }
  { \__flmx_default:nn { decofinal } {#1} }
\NewDocumentCommand \flmxDefineDecoMarginDefault { +m }
  { \__flmx_default:nn { decomargin } {#1} }
\NewDocumentCommand \flmxDefineDroppedDefault { +m }
  { \__flmx_default_ii:nn { dropped } {#1} }
\NewDocumentCommand \flmxDefineUnsupportedDefault { +m }
  { \__flmx_default_ii:nn { unsupported } {#1} }
\NewDocumentCommand \flmxDefineUnknownDefault { +m }
  { \__flmx_default_ii:nn { unknown } {#1} }

% --- the target of an anchor --------------------------------------------

\cs_new_protected:Npn \__flmx_target:n #1
  {
    \mode_if_math:TF
      {
        \legacy_if:nF { measuring@ }
          {
            \hypertarget { flmx:#1 } { }
            \protected@write \@auxout { }
              { \string \newlabel { flmx:#1 } { { } { \thepage } { } { flmx:#1 } { } } }
          }
      }
      { \phantomsection \label { flmx:#1 } }
  }
\cs_generate_variant:Nn \__flmx_target:n { e }
\NewDocumentCommand \flmxTarget { m } { \__flmx_target:e {#1} }

% --- blocks, spans, decorations --------------------------------------------

\NewDocumentEnvironment { flmxBlock } { O { } }
  { \__flmx_env_begin:nnn { block } { role, kinds } {#1} }
  { \__flmx_env_end:n { block } }
\flmxDefineBlockDefault { \par \__flmx_anchor: } { \par }

\NewDocumentCommand \flmxSpan { O { } +m }
  { \__flmx_command:nnnn { span } { role, kinds, style } {#1} {#2} }
\flmxDefineSpanDefault { \__flmx_anchor: #1 }
\flmxDefineSpan { emph } { \__flmx_anchor: \emph {#1} }
\flmxDefineSpan { bold } { \__flmx_anchor: \textbf {#1} }
\flmxDefineSpan { italic } { \__flmx_anchor: \textit {#1} }
\flmxDefineSpan { monospace } { \__flmx_anchor: \texttt {#1} }
\flmxDefineSpan { underline } { \__flmx_anchor: \underline {#1} }
\flmxDefineSpan { small-caps } { \__flmx_anchor: \textsc {#1} }
\flmxDefineSpan { superscript } { \__flmx_anchor: \textsuperscript {#1} }
\flmxDefineSpan { subscript } { \__flmx_anchor: \textsubscript {#1} }
\flmxDefineSpan { unresolved-reference }
  { \__flmx_anchor: \textcolor { flmxMarkColor } {#1} }

\NewDocumentCommand \flmxDecoRunin { O { } +m }
  { \__flmx_command:nnnn { decorunin } { } {#1} {#2} }
\flmxDefineDecoRuninDefault
  { \textbf { #1 \flmxKey { stop } } \hspace { \flmxRuninSep } }
\NewDocumentCommand \flmxDecoFinal { O { } +m }
  { \__flmx_command:nnnn { decofinal } { } {#1} {#2} }
\flmxDefineDecoFinalDefault { \leavevmode \unskip \nobreak \hfill \mbox {#1} }
\NewDocumentCommand \flmxDecoMargin { O { } +m }
  { \__flmx_command:nnnn { decomargin } { } {#1} {#2} }
\flmxDefineDecoMarginDefault { \marginpar { \footnotesize #1 } }

\ExplSyntaxOff
% flm.latex.def.flmx.xnotes
% flm.latex.def.flmx.xnotes -- notes: footnotes, endnotes and the other kinds.
%
% The xnotes feature writes a note by its kind's placement.
%
% A kind placed in a list (the default of the standard kind footnote): the
% note's mark in the running text, a superscript span of role
% flm.xnotes.xnote.mark holding a link of kind xnote to the note; the notes
% are a block of role flm.xnotes.<kind>.list, kind tokens xnotes and the
% kind's name, holding the kind's title as an unnumbered heading and a
% labeled list, each item's label a superscript span of role
% flm.xnotes.entry.mark:
%   Text.\flmxSpan[style=superscript, role=flm.xnotes.xnote.mark,
%     anchor=fnref-1]{\flmxLink[kind=xnote, href={\#fn-1}, target=fn-1,
%     doc-role=noteref, label={Footnote 1}]{1}}
%   ...
%   \begin{flmxBlock}[role=flm.xnotes.footnote.list, kinds={xnotes footnote},
%     doc-role=endnotes]
%   \flmxHeading[kind=xnotes, level=2, toc=false, plain={Footnotes}]{Footnotes}
%   \begin{flmxList}[kinds={numbered xnotes footnote labeled}, depth=1]
%   \flmxItem[label={\flmxSpan[style=superscript,
%     role=flm.xnotes.entry.mark]{1}}, anchor=fn-1]
%   The note.
%   \end{flmxList}
%   \end{flmxBlock}
%
% A kind placed in place: the note where it is written,
%   \flmxXnote[kind, anchor, doc-role]{<mark>}{<content>}
% with the keys
%   kind      the note's kind (footnote, or a kind the document configures);
%   anchor    the note's own anchor, where it has one;
%   doc-role  the document role, footnote;
% <mark> FLM's mark, a superscript span of role flm.xnotes.xnote.mark
% carrying the target of the mark's anchor, and <content> the note's
% content, which may hold several paragraphs.
%
% Defaults:
%   the block of role flm.xnotes.<kind>.list, registered under the kind
%     token xnotes, so for every kind: a block of paragraphs, the target of
%     its anchor first (the heading and the list keep their own looks);
%   the link kinds xnote (mark to note) and xnote-back (note back to its
%     mark) as the links chunk's default link;
%   \flmxXnote of every kind: a footnote at the bottom of the page carrying
%     FLM's mark, in the text and before the note (LaTeX's footnote counter
%     is not used), the note's text starting with the target of its anchor.
% The marks' spans are left to the span's default, the superscript of their
% style: registering one of their roles replaces that look.
%
% A template changes the look with
%   \flmxDefineBlock{xnotes}{<begin code>}{<end code>}  (every kind's list)
%   \flmxDefineBlock{flm.xnotes.<kind>.list}{<begin code>}{<end code>}
%   \flmxDefineSpan{flm.xnotes.xnote.mark}{<code over #1>}
%   \flmxDefineSpan{flm.xnotes.entry.mark}{<code over #1>}
%   \flmxDefineLink{xnote}{<code over #1>}, \flmxDefineLink{xnote-back}{...}
%   \flmxDefineXnote{<kind>}{<code over #1 (the mark) and #2 (the content)>}
% for instance \flmxDefineXnote{footnote}{#1\marginpar{\footnotesize #1 #2}}
% for margin notes; \flmxDefineXnoteDefault{...} changes every kind without
% a registration of its own.

\ExplSyntaxOn

\flmxDefineBlock { xnotes }
  { \par \flmxIfKey { anchor } { \flmxTarget { \flmxKey { anchor } } } { } }
  { \par }
\flmxDefineLink { xnote } { \__flmx_link:n {#1} }
\flmxDefineLink { xnote-back } { \__flmx_link:n {#1} }

\NewDocumentCommand \flmxXnote { O { } +m +m }
  { \__flmx_command:nnnnn { xnote } { kind } {#1} {#2} {#3} }

% The mark #1 as LaTeX's footnote mark shows it: the span of the mark's role
% set as it stands (\@makefnmark raises it), and, where #2 is false, without
% placing the mark's target, so that the target stands once, in the text.
\cs_new_protected:Npn \__flmx_xnote_mark:nn #1#2
  {
    \group_begin:
      \bool_if:nF {#2} { \cs_set_eq:NN \flmxTarget \use_none:n }
      \flmxDefineSpan { flm.xnotes.xnote.mark }
        { \flmxIfKey { anchor } { \flmxTarget { \flmxKey { anchor } } } { } ##1 }
      #1
    \group_end:
  }

% The two forms \@thefnmark holds: the mark placing its target (in the text)
% and the mark placing none (before the note, and where a \ref to the note's
% anchor typesets the mark). LaTeX internal names, not expl3 ones: LaTeX
% writes \@thefnmark to the .aux file as the label's text, which is read back
% outside expl3 syntax.
\cs_new_protected:Npn \flmx@xnotemark #1
  { \__flmx_xnote_mark:nn {#1} { \c_true_bool } }
\cs_new_protected:Npn \flmx@xnotemarknotarget #1
  { \__flmx_xnote_mark:nn {#1} { \c_false_bool } }

% Removes the paragraph breaks at the head of a note's content (a note of
% several paragraphs begins with one), so that its first line follows the
% mark.
\cs_new_protected:Npn \__flmx_xnote_skip_pars:
  { \peek_meaning_remove_ignore_spaces:NTF \par { \__flmx_xnote_skip_pars: } { } }

% The default: LaTeX's \@footnotemark and \@footnotetext, the pair \footnote
% calls, with the mark \@thefnmark set to FLM's mark.
\flmxDefineXnoteDefault
  {
    \cs_gset:Npn \@thefnmark { \flmx@xnotemark {#1} }
    \@footnotemark
    \cs_gset:Npn \@thefnmark { \flmx@xnotemarknotarget {#1} }
    \@footnotetext
      {
        \flmxIfKey { anchor } { \flmxTarget { \flmxKey { anchor } } } { }
        \__flmx_xnote_skip_pars: #2
      }
  }

\ExplSyntaxOff
% flm.latex.def.flmx.headings
% flm.latex.def.flmx.headings -- headings.
%
% FLM numbers the headings itself: the number stands in the title, in a span
% of role flm.heading.number, and no LaTeX counter is stepped (\section and
% secnumdepth play no part). The two elements:
%   \flmxHeading[kind, level, toc, anchor, number, plain, alt]{<title>}
%       a heading on a line of its own;
%   \flmxRuninHeading[kind, level, toc, anchor, number, plain, alt, stop]{<title>}
%       the title of a heading set at the start of its paragraph, inside
%       \flmxDecoRunin, which writes the title, the stop and \flmxRuninSep.
% Their keys:
%   kind     the heading kind's name (section, subsection, ...);
%   level    the level, counted from zero (section is 1);
%   toc      true or false: whether the heading is listed in a table of
%            contents;
%   anchor   where the heading is a reference target;
%   number   the heading's number as plain text, where the document numbered
%            it (it also stands in the title);
%   plain    the title's characters as plain text, the number included;
%   alt      the alternative title, where the author wrote one;
%   stop     (run-in) the punctuation after the title, where one is written.
%
% The default look. \flmxHeadingLevel{<level>}{<title>} typesets <title> as a
% heading of <level>, through LaTeX's \@startsection in its starred form with
% this table: level 0 the article's \section spacing, a page-break penalty
% before it and \huge; levels 1 to 5 the article's \section, \subsection,
% \subsubsection, \paragraph and \subparagraph (levels 4 and 5 run into the
% paragraph after them); a level past 5 as level 5. Inside the heading,
% before the title, it places the target of the anchor (\flmxTarget), a PDF
% bookmark where toc=true (\pdfbookmark at the level key itself, so that
% hyperref's bookmarksdepth counts FLM's levels as an article's: section 1,
% subsection 2; its text the plain key, its name the anchor, or a counter of
% this chunk where there is none), and, for levels 0 and 1, the running
% head: \markboth (level 0) or \markright (level 1) with the alt key where
% present, else plain. The target, the bookmark and the running head act
% once, when the heading itself is typeset: a \nameref to the heading
% typesets its title alone. The heading's keys must be in scope: call it
% from a heading's registered code.
%
% Dispatch. A heading is written with the code registered under its kind,
% then under its level, then with the default (\flmxHeadingLevel at the
% heading's level). No kind is registered by this chunk: a template changes
% the look of a kind by registering it (looked up first),
%   \flmxDefineHeading{section}{<code over #1, the title>}
% or the look of every heading of a level by registering the level,
%   \flmxDefineHeading{2}{<code over #1, the title>}
% \flmxRuninHeading's default sets the name a \nameref to the heading
% prints to the plain key, places the target, the bookmark and the running
% head, and writes the title; register a look with
% \flmxDefineRuninHeading{<kind or level>}{...}.

\ExplSyntaxOn

\int_new:N \g__flmx_bookmark_int
\tl_new:N \l__flmx_heading_name_tl
\tl_new:N \l__flmx_heading_places_tl
\tl_new:N \g__flmx_heading_places_tl
\cs_generate_variant:Nn \tl_set:Nn { Ne }
\cs_generate_variant:Nn \tl_put_right:Nn { Ne }

% Sets \l__flmx_heading_places_tl to the code placing the target, the
% bookmark and the running head of the heading whose keys are in scope, the
% keys' values written into it: a run-in heading of \@startsection is set
% at the start of the next paragraph, once the keys are out of scope.
\cs_new_protected:Npn \__flmx_heading_places:
  {
    \tl_clear:N \l__flmx_heading_places_tl
    \flmxIfKey { anchor }
      {
        \tl_put_right:Ne \l__flmx_heading_places_tl
          { \exp_not:N \flmxTarget { \exp_not:v { l__flmx_key_anchor_tl } } }
      }
      { }
    \str_if_eq:eeT { \flmxKey { toc } } { true }
      {
        \flmxIfKey { anchor }
          { \tl_set:Ne \l__flmx_heading_name_tl { flmxbookmark:\flmxKey { anchor } } }
          {
            \int_gincr:N \g__flmx_bookmark_int
            \tl_set:Ne \l__flmx_heading_name_tl
              { flmxbookmark.\int_use:N \g__flmx_bookmark_int }
          }
        \tl_put_right:Ne \l__flmx_heading_places_tl
          {
            \exp_not:N \pdfbookmark [ \flmxKey { level } ]
              { \flmxIfKey { plain } { \exp_not:v { l__flmx_key_plain_tl } } { } }
              { \exp_not:V \l__flmx_heading_name_tl }
          }
      }
    \int_case:nn { \flmxKey { level } }
      {
        { 0 }
          {
            \tl_put_right:Ne \l__flmx_heading_places_tl
              { \exp_not:N \markboth { \__flmx_heading_mark: } { } }
          }
        { 1 }
          {
            \tl_put_right:Ne \l__flmx_heading_places_tl
              { \exp_not:N \markright { \__flmx_heading_mark: } }
          }
      }
  }
% The running head's text, not expanded further: alt, else plain.
\cs_new:Npn \__flmx_heading_mark:
  {
    \flmxIfKey { alt }
      { \exp_not:v { l__flmx_key_alt_tl } }
      { \flmxIfKey { plain } { \exp_not:v { l__flmx_key_plain_tl } } { } }
  }

% The look of the table's level #1 over the title #2.
\cs_new_protected:Npn \__flmx_heading_level:nn #1#2
  {
    \__flmx_heading_places:
    \tl_gset_eq:NN \g__flmx_heading_places_tl \l__flmx_heading_places_tl
    \__flmx_heading_table:nn { \int_min:nn {#1} { 5 } } {#2}
  }
% Places the target, the bookmark and the running head of the heading being
% typeset, then forgets them. It stands before the title inside the title
% argument of \@startsection, which nameref stores as the heading's name: a
% \nameref typesetting that name meets it again and places nothing. Its name
% is a LaTeX internal one (with @), since it is written to the .aux file.
\cs_new_protected:Npn \flmx@headingplace
  {
    \tl_use:N \g__flmx_heading_places_tl
    \tl_gclear:N \g__flmx_heading_places_tl
  }
% #1 the level of the table, #2 the title.
\cs_new_protected:Npn \__flmx_heading_table:nn #1#2
  {
    \int_case:nn {#1}
      {
        { 0 }
          {
            \par \addpenalty { -\@highpenalty }
            \@startsection { flmxheading } { 0 } { \z@ }
              { -3.5ex \@plus -1ex \@minus -.2ex } { 2.3ex \@plus .2ex }
              { \normalfont \huge \bfseries } * { \flmx@headingplace #2 }
          }
        { 1 }
          {
            \@startsection { flmxheading } { 1 } { \z@ }
              { -3.5ex \@plus -1ex \@minus -.2ex } { 2.3ex \@plus .2ex }
              { \normalfont \Large \bfseries } * { \flmx@headingplace #2 }
          }
        { 2 }
          {
            \@startsection { flmxheading } { 2 } { \z@ }
              { -3.25ex \@plus -1ex \@minus -.2ex } { 1.5ex \@plus .2ex }
              { \normalfont \large \bfseries } * { \flmx@headingplace #2 }
          }
        { 3 }
          {
            \@startsection { flmxheading } { 3 } { \z@ }
              { -3.25ex \@plus -1ex \@minus -.2ex } { 1.5ex \@plus .2ex }
              { \normalfont \normalsize \bfseries } * { \flmx@headingplace #2 }
          }
        { 4 }
          {
            \@startsection { flmxheading } { 4 } { \z@ }
              { 3.25ex \@plus 1ex \@minus .2ex } { -1em }
              { \normalfont \normalsize \bfseries } * { \flmx@headingplace #2 }
          }
        { 5 }
          {
            \@startsection { flmxheading } { 5 } { \parindent }
              { 3.25ex \@plus 1ex \@minus .2ex } { -1em }
              { \normalfont \normalsize \bfseries } * { \flmx@headingplace #2 }
          }
      }
  }

\NewDocumentCommand \flmxHeadingLevel { m +m }
  { \__flmx_heading_level:nn {#1} {#2} }

\NewDocumentCommand \flmxHeading { O { } +m }
  {
    \group_begin:
      \__flmx_set_keys:n {#1}
      \__flmx_find:nn { heading } { kind, level }
      \use:c { \l__flmx_found_tl } {#2}
      % Keep the keys and formatting local, but carry LaTeX's continuation
      % into the scope of the following paragraph. For a displayed heading,
      % \@afterheading's \everypar clears the global \if@nobreak and applies
      % \if@afterindent; losing it suppresses later headings' before-space.
      % A run-in heading also needs its deferred title, \@svsechd. Preserve
      % the name for a following \label in either case. These assignments
      % are local to the containing scope, not global across environments.
      % Do this at dispatch so registered \@startsection code works too.
      \use:e
        {
          \group_end:
          \exp_not:N \everypar { \exp_not:V \everypar }
          \legacy_if:nTF { @afterindent }
            { \exp_not:N \@afterindenttrue }
            { \exp_not:N \@afterindentfalse }
          \exp_not:N \cs_set:Npn \exp_not:N \@currentlabelname
            { \exp_not:V \@currentlabelname }
          \legacy_if:nT { @noskipsec }
            {
              \exp_not:N \cs_set:Npn \exp_not:N \@svsechd
                { \exp_not:V \@svsechd }
            }
        }
  }
\flmxDefineHeadingDefault { \flmxHeadingLevel { \flmxKey { level } } {#1} }

\NewDocumentCommand \flmxRuninHeading { O { } +m }
  { \__flmx_command:nnnn { runinheading } { kind, level } {#1} {#2} }
\flmxDefineRuninHeadingDefault
  {
    \flmxIfKey { plain }
      { \tl_set:Nv \@currentlabelname { l__flmx_key_plain_tl } }
      { }
    \__flmx_heading_places: \tl_use:N \l__flmx_heading_places_tl #1
  }

\ExplSyntaxOff
% flm.latex.def.flmx.links
% flm.latex.def.flmx.links -- links.
%
%   \flmxLink[kind, href, target, doc-role, label, shows-target]{<content>}
%       a link around <content>. Its keys:
%         kind          what the link is: ref (a reference), term, cite,
%                       xnote, xnote-back, toc, url, or a kind a feature adds;
%         href          the address, always written (#<anchor> for a link
%                       inside the document);
%         target        for a link inside the document, the anchor the link
%                       leads to, without the #: \hyperref[flmx:<target>]{...}
%                       and \pageref{flmx:<target>} reach it;
%         doc-role      the document role of the link, where it has one;
%         label         a text naming the link for assistive technology;
%         shows-target  true where the content shows the address itself.
%
% The default of every kind: a link with a target is
% \hyperref[flmx:<target>]{<content>}, any other \href{<href>}{<content>};
% where shows-target=true, the content is set in \flmxUrlFont. The kinds ref
% and url are registered here with that code; the kinds the features write
% are registered by their definition pieces (toc by the toc feature's, with
% leaders and the page number; term by the defterm feature's; cite by the
% citations feature's; xnote and xnote-back by the xnotes feature's), so that
% a template changes one kind with
%   \flmxDefineLink{<kind>}{<code over #1, the content>}
% and every kind with \flmxDefineLinkDefault{...} and the registrations.
% \flmxUrlFont is the font of a link that shows its address, \ttfamily;
% change it with \RenewDocumentCommand\flmxUrlFont{}{...}.

\ExplSyntaxOn
\NewDocumentCommand \flmxUrlFont { } { \ttfamily }

\cs_new_protected:Npn \__flmx_link:n #1
  {
    \flmxIfKey { target }
      { \hyperref [ flmx:\flmxKey { target } ] { \__flmx_link_content:n {#1} } }
      { \exp_args:Ne \href { \flmxKey { href } } { \__flmx_link_content:n {#1} } }
  }
\cs_new_protected:Npn \__flmx_link_content:n #1
  {
    \str_if_eq:eeTF { \flmxKey { shows-target } } { true }
      { { \flmxUrlFont #1 } }
      {#1}
  }

\NewDocumentCommand \flmxLink { O { } +m }
  { \__flmx_command:nnnn { link } { kind } {#1} {#2} }
\flmxDefineLinkDefault { \__flmx_link:n {#1} }
\clist_map_inline:nn { ref, url }
  { \flmxDefineLink {#1} { \__flmx_link:n {##1} } }
\ExplSyntaxOff
% flm.latex.def.flmx.lists
% flm.latex.def.flmx.lists -- lists.
%
% FLM computes every label itself (a number, a bullet, a term): the label
% stands in the item's keys, and no LaTeX counter is stepped. The elements:
%   \begin{flmxList}[kinds, anchor, depth]
%   \flmxItem[label, anchor]
%   <the item's content: paragraphs, verbatim blocks, nested lists>
%   \flmxItem[...]
%   ...
%   \end{flmxList}
% \flmxItem takes no content argument: the item's content follows it, up to
% the next \flmxItem or the \end, as with LaTeX's \item. Their keys:
%   kinds    (list) the list's kind tokens: at most one of numbered,
%            unnumbered, terms first, then others, such as the name of the
%            environment the list comes from (enumerate, itemize,
%            description) and labeled (every item carries its own label);
%   anchor   (list, item) where the list or the item is a reference target;
%   depth    (list) the list's depth among lists, 1 for an outermost list;
%   label    (item) the item's label (its number, its bullet, its term),
%            where it has one.
%
% The default look, over LaTeX's \list and \item, by the list's structural
% kind tokens (the only ones it knows):
%   numbered, labeled  \item[<label>], the label flush right in the label
%            box; an item of a labeled list without a label shows none;
%   unnumbered  \item[<label>] where the item has a label, else \item with a
%            bullet by depth: \textbullet, \textendash,
%            \textasteriskcentered, then \textperiodcentered;
%   terms    the label in bold, run into the item (\labelwidth zero), the
%            item's lines hanging at the list's indentation.
% A list without a structural token is written as unnumbered. The list's
% target stands before the list, an item's at the start of its content.
%
% Registering a look. A list and its items are written with the code
% registered under the first of the list's kind tokens that has one, else
% with the default. Both registrations go by the list's kind tokens (an
% item has none of its own):
%   \flmxDefineList{<kind token>}{<begin code>}{<end code>}
%   \flmxDefineItem{<kind token>}{<code run at the item>}
% for instance \flmxDefineList{enumerate}{\begin{enumerate}}{\end{enumerate}}
% and \flmxDefineItem{enumerate}{\item[\flmxKey{label}]}. An item's code
% sees the item's keys (\flmxKey{label}); the list's kind tokens and depth
% are \flmxListKinds and \flmxListDepth. The defaults are replaced with
% \flmxDefineListDefault{<begin>}{<end>} and \flmxDefineItemDefault{<code>}.

\ExplSyntaxOn

\tl_new:N \l__flmx_list_kinds_tl
\tl_new:N \l__flmx_list_depth_tl
\tl_new:N \l__flmx_list_shape_tl
\bool_new:N \l__flmx_list_labeled_bool
\cs_generate_variant:Nn \str_case:nnF { V }

\NewExpandableDocumentCommand \flmxListKinds { } { \tl_use:N \l__flmx_list_kinds_tl }
\NewExpandableDocumentCommand \flmxListDepth { } { \tl_use:N \l__flmx_list_depth_tl }

% Saves the list's kind tokens and depth, and its shape: the structural
% token the default goes by (unnumbered where there is none), and whether
% it is labeled.
\cs_new_protected:Npn \__flmx_list_save:
  {
    \tl_set:Ne \l__flmx_list_kinds_tl { \flmxKey { kinds } }
    \tl_set:Ne \l__flmx_list_depth_tl { \flmxKey { depth } }
    \tl_set:Nn \l__flmx_list_shape_tl { unnumbered }
    \bool_set_false:N \l__flmx_list_labeled_bool
    \tl_set_eq:NN \l__flmx_kinds_tl \l__flmx_list_kinds_tl
    \tl_replace_all:Nnn \l__flmx_kinds_tl { ~ } { , }
    \clist_map_inline:Nn \l__flmx_kinds_tl
      {
        \str_case:nn {##1}
          {
            { numbered } { \tl_set:Nn \l__flmx_list_shape_tl { numbered } }
            { terms } { \tl_set:Nn \l__flmx_list_shape_tl { terms } }
            { labeled } { \bool_set_true:N \l__flmx_list_labeled_bool }
          }
      }
  }

\NewDocumentEnvironment { flmxList } { O { } }
  {
    \__flmx_set_keys:n {#1}
    \__flmx_list_save:
    \__flmx_env_begin:nnn { list } { kinds } {#1}
  }
  { \__flmx_env_end:n { list } }

% The bullet of an unnumbered list at the depth #1.
\cs_new:Npn \__flmx_list_bullet:n #1
  {
    \int_case:nnF {#1}
      {
        { 1 } { \textbullet }
        { 2 } { \textendash }
        { 3 } { \textasteriskcentered }
      }
      { \textperiodcentered }
  }
\cs_new_protected:Npn \__flmx_list_term:n #1
  { \hspace \labelsep \normalfont \bfseries #1 }

\flmxDefineListDefault
  {
    \__flmx_anchor:
    \str_case:VnF \l__flmx_list_shape_tl
      {
        { terms }
          {
            \list { }
              {
                \setlength \labelwidth { 0pt }
                \setlength \itemindent { -\leftmargin }
                \cs_set_eq:NN \makelabel \__flmx_list_term:n
              }
          }
        { numbered } { \list { } { } }
      }
      { \list { \__flmx_list_bullet:n { \l__flmx_list_depth_tl } } { } }
  }
  { \endlist }

% The item: its keys, then the code registered under the list's first kind
% token that has one, else the default.
\cs_new_protected:Npn \__flmx_item_find:
  {
    \bool_set_false:N \l__flmx_found_bool
    \tl_set:Nn \l__flmx_found_tl { __flmx_item_default }
    \tl_set_eq:NN \l__flmx_kinds_tl \l__flmx_list_kinds_tl
    \tl_replace_all:Nnn \l__flmx_kinds_tl { ~ } { , }
    \clist_map_inline:Nn \l__flmx_kinds_tl { \__flmx_try:nn { item } {##1} }
  }
\NewDocumentCommand \flmxItem { O { } }
  {
    \__flmx_set_keys:n {#1}
    \__flmx_item_find:
    \use:c { \l__flmx_found_tl }
  }
\flmxDefineItemDefault
  {
    \flmxIfKey { label }
      { \item [ \flmxKey { label } ] }
      { \bool_if:NTF \l__flmx_list_labeled_bool { \item [ ] } { \item } }
    \__flmx_anchor:
  }

\ExplSyntaxOff
\endinput
